\section{Patch positivity and the critical profile}
\label{sec:patch-positivity}

The explicit patch factors separate two positivity questions.  Bulk
positivity depends only on the flip-rate coefficients.  Once it holds,
the remaining obstruction is an affine threshold in the terminal
one-density.

\subsection{Bulk positivity}

\begin{definition}
The spin system has the \emph{patch positivity property} when every
bulk contribution \(C(P)\) is nonnegative.
\end{definition}

\begin{theorem}[Local criterion]
\label{thm:patch-positivity-criterion}
All bulk contributions based at \(i\) are nonnegative if and only if,
for every nonempty \(S\subseteq N(i)\),
\begin{equation}
  \begin{cases}
  c_i^0(S)+c_i^1(S)\le0,\\
  c_i^1(\vn)c_i^0(S)-c_i^0(\vn)c_i^1(S)\ge0,
  \end{cases}
  \qquad
  \text{if }c_i^0(\vn)+c_i^1(\vn)>0,
  \label{eq:patch-positive-nondegenerate}
\end{equation}
whereas in the degenerate case
\(c_i^0(\vn)+c_i^1(\vn)=0\) the criterion is
\begin{equation}
  c_i^0(S)+c_i^1(S)\le0,
  \qquad
  c_i^1(S)\le0.
  \label{eq:patch-positive-degenerate}
\end{equation}
\end{theorem}

\begin{proof}
The four bulk rows in \cref{prop:explicit-patch-factors} can be checked
separately.  The \(\sI\sI\) and \(\sI\sO\) rows are automatically
nonnegative.  The \(\sO\sO\) row is nonnegative exactly when
\(c_i^0(S)+c_i^1(S)\le0\).  Under that condition, the numerator in the
\(\sO\sI\) row is smallest at the limiting empty-neighbour density.
When the empty-neighbour refresh rate is positive, its nonnegativity
is exactly the determinant inequality in
\eqref{eq:patch-positive-nondegenerate}.  When it vanishes, the same
numerator reduces to \(-c_i^1(S)\).  The full calculation is given in
\cref{app:patch-calculations}.
\end{proof}

\subsection{Critical density}

\begin{definition}
For a patch-positive system, the patch critical density at \(i\) is
\[
  p_i^\star
  =
  \inf\set{
    p\in[0,1]:
    C(p,P)\ge0
    \text{ for every end patch based at }i
  }.
\]
The resulting profile is \(\pstar=(p_i^\star)_{i\in\Lambda}\).
\end{definition}

\begin{proposition}[Coefficient formula]
\label{prop:critical-density-formula}
For every \(i\in\Lambda\),
\begin{equation}
  p_i^\star
  =
  \max\left\{
    0,
    \sup_{\substack{\vn\ne S\subseteq N(i)\\
                    c_i^0(S)+c_i^1(S)<0}}
    \frac{c_i^0(S)}{c_i^0(S)+c_i^1(S)}
  \right\}.
  \label{eq:critical-density-formula}
\end{equation}
If
\[
  r_i=c_i^0(\vn)+c_i^1(\vn)>0,
  \qquad
  q_i=\frac{c_i^0(\vn)}{r_i},
\]
then \(p_i^\star\le q_i\).
\end{proposition}

\begin{proof}
The \(\sI\sE\) end factor is nonnegative for every terminal value.
For an \(\sO\sE\) patch with initial target \(S\), the numerator is
\[
  c_i^0(S)
  -
  \bigl(c_i^0(S)+c_i^1(S)\bigr)\psi_i(\Delta,p).
\]
Patch positivity makes this expression nondecreasing in \(p\).  As
the patch length varies, its potentially smallest value is the
zero-length value
\[
  c_i^0(S)
  -
  \bigl(c_i^0(S)+c_i^1(S)\bigr)p.
\]
Solving the resulting inequalities over all nonempty \(S\) gives
\eqref{eq:critical-density-formula}.  The determinant inequality in
\eqref{eq:patch-positive-nondegenerate} gives \(p_i^\star\le q_i\).
\end{proof}

\subsection{Centered classes}

For \(A\Subset\Lambda\), define
\begin{equation}
  \chi_A^\star(\eta)
  =
  \prod_{i\in A}\bigl(\eta(i)-p_i^\star\bigr),
  \qquad
  \chi_\vn^\star=1.
  \label{eq:centered-monomial}
\end{equation}

\begin{definition}
The high-density class and its lower extension are
\begin{align}
  \Mstar
  &=
  \set{
    \mu:
    \mu(\chi_A^\star)\ge0
    \text{ for every }A\Subset\Lambda
  },
  \label{eq:Mstar}
  \\
  \Mminus
  &=
  \set{
    \mu:
    \mu(\chi_A^\star)
    \ge-\chi_A^\star(\mathbf1)
    \text{ for every }A\Subset\Lambda
  }.
  \label{eq:Mminus}
\end{align}
\end{definition}

Both classes are convex and weakly closed, and
\(\Mstar\subseteq\Mminus\).  Moreover,
\begin{equation}
  \mu\in\Mminus
  \quad\Longleftrightarrow\quad
  \frac{\mu+\mu_{\mathbf1}}2\in\Mstar,
  \label{eq:Mminus-affine}
\end{equation}
where \(\mu_{\mathbf1}\) is the point mass at the all-one
configuration.

For a Bernoulli product law with one-density profile \(\mathbf p\),
\[
  \mu_{\mathbf p}(\chi_A^\star)
  =
  \prod_{i\in A}(p_i-p_i^\star).
\]
Consequently
\begin{align}
  \mu_{\mathbf p}\in\Mstar
  &\quad\Longleftrightarrow\quad
  \mathbf p\ge\pstar,
  \label{eq:product-Mstar}
  \\
  \mu_{\mathbf p}\in\Mminus
  &\quad\Longleftrightarrow\quad
  p_i\ge(2p_i^\star-1)\vee0
  \quad\text{for every }i.
  \label{eq:product-Mminus}
\end{align}
These product descriptions do not characterize the convex classes:
both also contain non-product laws and laws that are not mixtures of
the displayed product measures.

\begin{lemma}
\label{lem:all-measures-Mminus}
If \(\pstar\le\frac12\mathbf1\), then every probability measure on
\(\{0,1\}^{\Lambda}\) belongs to \(\Mminus\).
\end{lemma}

\begin{proof}
For every configuration \(\eta\),
\[
  \frac{\chi_A^\star(\eta)}
       {\chi_A^\star(\mathbf1)}
  =
  \prod_{i\in A}
  \frac{\eta(i)-p_i^\star}{1-p_i^\star}
  \ge-1.
\]
Integrating proves the defining inequality
\eqref{eq:Mminus}.
\end{proof}
